\chapter{Funções e Variações}\label{ch:g11:func}

Este capítulo organiza o que é preciso saber sobre funções antes
cálculo: a pequena família de \emph{\omterm{def:g11:func:reference}{funções de referência}} cujo \omterm{def:g10:functions:graph}{gráficos} deve
ser conhecido de cor, o significado preciso de \emph{crescente} e
\emph{decrescente}e como as variações se comportam quando funções são combinados.

\section{Funções e curvas}

\begin{definition}[Função, domínio, curva]\label{def:g11:func:function}
UM \emph{função}\index{função} $f $ atribui a cada número $x$ de um conjunto
$ D \subseteq \R $ (isso é \emph{domínio domínio}\index{domínio domínio}) exatamente um número
$ f(x)$. Isso é \emph{curva} (ou grafo) é o conjunto de pontos
$\bigl(x, f(x)\bigr)$ para $ x \in D $.
\end{definition}

\begin{example}
A fórmula $ f(x) = \frac{\sqrt{x+2}}{x-1}$ requer $ x + 2 \geq 0$ (para
a \omterm{def:g11:quad:discriminant}{raiz} quadrada) e $ x \neq 1$ (denominador diferente de zero): o domínio é
$\intco{-2}{1} \cup \intoo{1}{+\infty}$.
\end{example}

\section{As funções de referência}

Os cinco seguintes funções e suas curvas devem tornar-se reflexos.

\begin{definition}[Funções de referência]\label{def:g11:func:reference}
\index{funções de referência}
\[
x \mapsto x^2 \text{ (square)}, \quad
x \mapsto x^3 \text{ (cube)}, \quad
x \mapsto \sqrt{x}, \quad
x \mapsto \frac{1}{x}, \quad
x \mapsto \abs{x}.
\]
Deles domínios são $\R $, $\R $, $\intco{0}{+\infty}$, $\R \setminus \{0\}$
e $\R $ respectivamente. A curva de $ x \mapsto \frac1x $ é um
\emph{hipérbole}\index{hipérbole}.
\end{definition}

\begin{omfigure}
\begin{tikzpicture}
\begin{axis}[omaxis,
  width=7cm, height=5.6cm,
  xmin=-2.3, xmax=2.3, ymin=-3.4, ymax=3.6,
  xtick={-1,1}, ytick={-1,1}]
  \addplot[omDef, thick, domain=-1.85:1.85] {x^2};
  \addplot[omThm, thick, domain=-1.5:1.5] {x^3};
  \node[omDef, font=\small, anchor=east] at (axis cs:-1.5,1.9) {$ x^2$};
  \node[omThm, font=\small, anchor=east] at (axis cs:-1.32,-1.73) {$ x^3$};
\end{axis}
\end{tikzpicture}%
\begin{tikzpicture}
\begin{axis}[omaxis,
  width=7cm, height=5.6cm,
  xmin=-3.4, xmax=4.4, ymin=-3.4, ymax=3.6,
  xtick={1}, ytick={1}]
  \addplot[omProp, thick, domain=-3.2:-0.31, samples=60] {1/x};
  \addplot[omProp, thick, domain=0.31:4.2, samples=60] {1/x};
  \addplot[omDef, thick, domain=0:4.2] {sqrt(x)};
  \addplot[omThm, thick, domain=-3.2:3.2] {abs(x)};
  \node[omDef, font=\small, anchor=north] at (axis cs:3.6,1.55)
    {$\sqrt{x}$};
  \node[omThm, font=\small, anchor=east] at (axis cs:-2.4,2.8)
    {$\abs{x}$};
  \node[omProp, font=\small, anchor=west] at (axis cs:0.75,2.4)
    {$\tfrac1x $};
\end{axis}
\end{tikzpicture}

{\small As cinco curvas de referência: a \omterm{def:g11:quad:parabola}{parábola} $ x^2$ e o cubo $ x^3$
(esquerda); $\sqrt x $, $\abs x $ e o \omterm{def:g11:func:reference}{hipérbole} $\frac1x $ (certo).}
\end{omfigure}

\section{Variações em um intervalo}

\begin{definition}[Aumentando, diminuindo]\label{def:g11:func:monotone}
Seja $f$ definida em um \omterm{def:g10:numbers:interval}{intervalo} $ I $. A \omterm{def:g11:func:function}{função} $f$ é
\emph{aumentando}\index{aumentando} sobre $ I $ se preserva a ordem:
\[
\text{para todo } u, v \in I: \quad u < v \implies f(u) \leq f(v),
\]
e \emph{decrescente} se inverter a ordem ($ u < v \implies f(u) \geq
f(v)$). Com desigualdades estritas à direita, $f$ é \emph{estritamente}
aumentando (resp.\ decrescente). UM \omterm{def:g11:func:function}{função} que está aumentando ou
decrescente sobre $ I $ é \emph{monótono} sobre $ I $.
\end{definition}

\begin{method}[Provando uma variação da definição]\label{met:g11:func:variation}
Pegue dois números arbitrários $ u < v $ no \omterm{def:g10:numbers:interval}{intervalo}, e determine o sinal
da diferença $ f(v) - f(u)$, geralmente por \emph{\omterm{def:g10:algebra:expand}{factoring}} isto. Se o
sinal é constante no \omterm{def:g10:numbers:interval}{intervalo}, conclua.
\end{method}

\begin{proposition}[Variações das funções de referência]\label{prop:g11:func:refvariations}
\begin{enumerate}
  \item $ x^2$ é estritamente decrescente sobre $\intoc{-\infty}{0}$ e estritamente
        crescente sobre $\intco{0}{+\infty}$.
  \item $ x^3$ é estritamente crescente sobre $\R $.
  \item $\sqrt x $ é estritamente crescente sobre $\intco{0}{+\infty}$.
  \item $\frac1x $ é estritamente decrescente sobre $\intoo{-\infty}{0}$ e assim por diante
        $\intoo{0}{+\infty}$ (mas \emph{não} em seu \omterm{def:g10:numbers:interunion}{união}).
  \item $\abs x $ é estritamente decrescente sobre $\intoc{-\infty}{0}$ e
        estritamente crescente sobre $\intco{0}{+\infty}$.
\end{enumerate}
\end{proposition}

\begin{proof}
Seja $ u < v $. Nós fatoramos $ f(v) - f(u)$ em cada caso.

\emph{1.} $ v^2 - u^2 = (v-u)(v+u)$. Sobre $\intco{0}{+\infty}$: $ v - u > 0$
e $ v + u > 0$ (como $ v > u \geq 0$), então $ v^2 > u^2$. Sobre
$\intoc{-\infty}{0}$: $ v + u < 0$, então $ v^2 < u^2$.

\emph{2.} $ v^3 - u^3 = (v - u)(v^2 + uv + u^2)$, e o segundo fator
é igual $\left(v + \frac u2\right)^2 + \frac{3u^2}{4} > 0$ a menos que
$ u = v = 0$. Por isso $ v^3 > u^3$ sempre que $ v > u $.

\emph{3.} Para $0 \leq u < v $:
$\sqrt v - \sqrt u = \frac{(\sqrt v - \sqrt u)(\sqrt v + \sqrt u)}
{\sqrt v + \sqrt u} = \frac{v - u}{\sqrt v + \sqrt u} > 0$.

\emph{4.} Para $0 < u < v $:
$\frac1v - \frac1u = \frac{u - v}{uv} < 0$ desde $ u - v < 0$ e $ uv > 0$;
mesmo cálculo para $ u < v < 0$. No \omterm{def:g10:numbers:interunion}{união} a reivindicação falha:
$-1 < 1$ mas $\frac{1}{-1} = -1 < 1 = \frac11$.

\emph{5.} Sobre $\intco{0}{+\infty}$, $\abs x = x $ é estritamente crescente; sobre
$\intoc{-\infty}{0}$, $\abs x = -x $ é estritamente decrescente.
\end{proof}

\begin{remark}
O ponto 4 é uma armadilha clássica: ``decrescente sobre $\intoo{-\infty}{0}$ e assim por diante
$\intoo{0}{+\infty}$'' é \emph{não} o mesmo que ``decrescente sobre
$\R \setminus \{0\}$''. Variações só fazem sentido em \omterm{def:g10:numbers:interval}{intervalos}.
\end{remark}

\begin{example}[Comparando sem calcular]
Desde $\sqrt x $ é estritamente crescente, $\sqrt{7} < \sqrt{8}$; desde
$ x^2$ é estritamente decrescente sobre $\intoc{-\infty}{0}$,
$(-3)^2 > (-2)^2$. A monotonicidade transforma comparações de imagens em
comparações de argumentos.
\end{example}

\section{Novas funções das antigas}

\begin{proposition}[Transformações]\label{prop:g11:func:transforms}
Seja $ u $ seja uma \omterm{def:g11:func:function}{função} definida em um \omterm{def:g10:numbers:interval}{intervalo} $ I $ e $ k \in \R $.
\begin{enumerate}
  \item $ u + k $ tem o \emph{\omterm{def:g11:func:parity}{par}} variações como $ u $ sobre $ I $; sua curva é
        o de $ u $ deslocado verticalmente por $ k $.
  \item Para $\lambda > 0$, $\lambda u $ tem as mesmas variações que $ u $; para
        $\lambda < 0$, as variações são invertidas.
  \item Se $ u \geq 0$ sobre $ I $, então $\sqrt u $ tem as mesmas variações que
        $ u $.
  \item Se $ u \neq 0$ e $ u $ tem login constante $ I $, então $\frac1u $ tem
        as variações \emph{oposto} para aqueles de $ u $.
\end{enumerate}
\end{proposition}

\begin{proof}
Seja $ u < v $ em $ I $ (abusando da notação, comparamos $x$-valores $ u < v $; escrever
$ s = u(x_1)$, $ t = u(x_2)$ para o imagens de $ x_1 < x_2$).

\emph{1.} $(u+k)(x_2) - (u+k)(x_1) = t - s $: a diferença de imagens é
inalterado, então seu sinal permanece inalterado.

\emph{2.} $\lambda t - \lambda s = \lambda(t - s)$: o sinal é mantido se
$\lambda > 0$, inverteu se $\lambda < 0$.

\emph{3.} Se $ u $ é crescente e $ x_1 < x_2$, então $0 \leq s \leq t $, e
$\sqrt{}$ ser crescente (\cref{prop:g11:func:refvariations}) dá
$\sqrt s \leq \sqrt t $. O mesmo argumento no decrescente caso.

\emph{4.} Se $ u $ é crescente, $0 < s \leq t $ (ou $ s \leq t < 0$), então
$\frac1s \geq \frac1t $ porque $ x \mapsto \frac1x $ é decrescente em cada
lado de $0$: a ordem é invertida.
\end{proof}

\begin{example}\label{ex:g11:func:composite}
Estudar $ f(x) = \sqrt{x^2 + 1}$ sobre $\R $. O interior \omterm{def:g11:func:function}{função}
$ u(x) = x^2 + 1$ é positivo, decrescente sobre $\intoc{-\infty}{0}$ e
crescente sobre $\intco{0}{+\infty}$ (um quadrado deslocado). Pelo ponto 3, $f$
tem as mesmas variações: decrescente então crescente, com \omterm{def:g10:functions:extrema}{mínimo}
$ f(0) = 1$.
\end{example}

\begin{omfigure}
\begin{tikzpicture}
\begin{axis}[omaxis,
  width=8.6cm, height=5.8cm,
  xmin=-2.6, xmax=2.6, ymin=-1.7, ymax=4.6,
  xtick={-1,1}, ytick={1,2}]
  \addplot[omDef, thick, domain=-1.9:1.9] {x^2};
  \addplot[omThm, thick, domain=-1.55:1.55] {x^2 + 1};
  \addplot[omProp, thick, domain=-2.4:2.4, samples=80]
    {sqrt(x^2 + 1)};
  \node[omDef, font=\small, anchor=west] at (axis cs:1.9,3.4) {$ x^2$};
  \node[omThm, font=\small, anchor=west] at (axis cs:1.35,3.3)
    {$ x^2{+}1$};
  \node[omProp, font=\small, anchor=north] at (axis cs:-1.9,1.95)
    {$\sqrt{x^2+1}$};
\end{axis}
\end{tikzpicture}

{\small Transformações em ação: mudança $ x^2$ por $1$ (vermelho) mantém
suas variações; tirar a \omterm{def:g11:quad:discriminant}{raiz} quadrada (laranja) os mantém novamente, como em
\cref{ex:g11:func:composite}.}
\end{omfigure}

\section{Funções pares e ímpares}

\begin{definition}[Paridade]\label{def:g11:func:parity}
Seja $f$ definida em um domínio $ D $ simétrico sobre $0$ (ou seja,\ $-x \in D $
sempre que $ x \in D $). A \omterm{def:g11:func:function}{função} $f$ é \emph{par}\index{função par}
se $ f(-x) = f(x)$ para todos $ x \in D $, e \emph{ímpar}\index{função ímpar} se
$ f(-x) = -f(x)$ para todos $ x \in D $.
\end{definition}

\begin{proposition}[Significado geométrico]\label{prop:g11:func:paritysym}
A curva de uma \omterm{def:g11:func:parity}{função par} é simétrica em relação ao $ y $-eixo; a curva
de uma \omterm{def:g11:func:parity}{função ímpar} é simétrica em relação à origem.
\end{proposition}

\begin{proof}
Se $f$ é \omterm{def:g11:func:parity}{par} e $(x, y)$ está na curva (então $ y = f(x)$), então
$ f(-x) = f(x) = y $: o ponto do espelho $(-x, y)$ também está na curva. Se
$f$ é \omterm{def:g11:func:parity}{ímpar}, $ f(-x) = -y $: o ponto $(-x, -y)$, simétrico de $(x,y)$ sobre
a origem, está na curva.
\end{proof}

\begin{example}
$ x^2$ e $\abs x $ são \omterm{def:g11:func:parity}{pares}; $ x^3$ e $\frac1x $ são \omterm{def:g11:func:parity}{ímpares};
$ f(x) = x^2 + x $ não é nenhum dos dois: $ f(-1) = 0$ enquanto $ f(1) = 2$, então
$ f(-1) \neq \pm f(1)$. A paridade reduz o trabalho pela metade: uma \omterm{def:g11:func:parity}{par} ou \omterm{def:g11:func:parity}{função ímpar}
só precisa ser estudado $\intco{0}{+\infty}$.
\end{example}

\section{Exercícios}

\begin{exercise}[$\star $]\label{exo:g11:func:1}
Dê o domínio de cada \omterm{def:g11:func:function}{função}:
\[
f(x) = \sqrt{3 - x}, \qquad
g(x) = \frac{1}{x^2 - 4}, \qquad
h(x) = \frac{\sqrt{x}}{x - 5} .
\]
\end{exercise}

\begin{exercise}[$\star $]\label{exo:g11:func:2}
Usando as variações do \omterm{def:g11:func:reference}{funções de referência}, compare sem
calculadora:
\[
\sqrt{11} \text{ e } \sqrt{13}; \qquad
(-2.1)^2 \text{ e } (-2.05)^2; \qquad
\frac{1}{0.9} \text{ e } \frac{1}{1.1}; \qquad
(-1.01)^3 \text{ e } (-0.99)^3 .
\]
\end{exercise}

\begin{exercise}[$\star $]\label{exo:g11:func:3}
Determine se cada \omterm{def:g11:func:function}{função} é \omterm{def:g11:func:parity}{par}, \omterm{def:g11:func:parity}{ímpar}, ou nenhum:
\[
f(x) = 3x^4 - x^2, \qquad
g(x) = x^3 + x, \qquad
h(x) = x^2 + x^3, \qquad
k(x) = \frac{x}{x^2+1} .
\]
\end{exercise}

\begin{exercise}[$\star $]\label{exo:g11:func:4}
UM \omterm{def:g11:func:function}{função} $f$ tem o seguinte comportamento de variação: crescente sobre
$\intcc{-3}{1}$ de $ f(-3) = -2$ a $ f(1) = 4$, então decrescente sobre
$\intcc{1}{5}$ até $ f(5) = 0$. Quantas soluções o equação
$ f(x) = 1$ usar $\intcc{-3}{5}$? E $ f(x) = 5$?
\end{exercise}

\begin{exercise}[$\star\star $]\label{exo:g11:func:5}
Usando \cref{met:g11:func:variation} (sinal de $ f(v) - f(u)$), prove que
$ f(x) = x^2 - 4x $ é estritamente crescente sobre $\intco{2}{+\infty}$.
\end{exercise}

\begin{exercise}[$\star\star $]\label{exo:g11:func:6}
Sem calcular nenhuma derivada, forneça as variações de
\[
f(x) = \sqrt{5 - x} \ \text{ on } \intoc{-\infty}{5},
\quad
g(x) = \frac{1}{x^2 + 1} \ \text{ on } \intco{0}{+\infty},
\quad
h(x) = 3 - 2\sqrt{x} \ \text{ on } \intco{0}{+\infty}.
\]
\end{exercise}

\begin{exercise}[$\star\star $]\label{exo:g11:func:7}
Seja $ f(x) = (x-2)^2 - 3$ sobre $\R $.
\begin{enumerate}
  \item Dê as variações de $f$ e seu \omterm{def:g10:functions:extrema}{mínimo}.
  \item Deduza as variações de $ g(x) = \frac{1}{(x-2)^2+1}$ e seu
        \omterm{def:g10:functions:extrema}{máximo}. (Observação $ g = \frac{1}{f + 4}$.)
\end{enumerate}
\end{exercise}

\begin{exercise}[$\star\star $]\label{exo:g11:func:8}
A curva de uma \omterm{def:g11:func:function}{função} $f$ passa pelos pontos $(0, 1)$, $(2, 3)$
e $(4, 1)$, e $f$ é crescente sobre $\intcc{0}{2}$ e decrescente sobre
$\intcc{2}{4}$. Esboce uma curva possível e depois esboce as curvas de
$ f + 2$, de $-f $, e de $2f $.
\end{exercise}

\begin{exercise}[$\star\star $]\label{exo:g11:func:9}
Mostre que a \omterm{def:g11:func:function}{função} $ f(x) = \frac{2x+1}{x+1}$ pode ser escrito
$ f(x) = 2 - \frac{1}{x+1}$, e deduza suas variações em
$\intoo{-1}{+\infty}$.
\end{exercise}

\begin{exercise}[$\star\star $]\label{exo:g11:func:10}
Verdadeiro ou falso (justifique ou dê um contraexemplo):
\begin{enumerate}
  \item A soma de dois funções crescentes sobre $ I $ é crescente sobre $ I $.
  \item O produto de dois funções crescentes sobre $ I $ é crescente sobre
        $ I $.
  \item Se $f$ é crescente sobre $ I $, então $ f^2$ ($= f \times f $) é
        crescente sobre $ I $.
\end{enumerate}
\end{exercise}

\begin{exercise}[$\star\star\star $]\label{exo:g11:func:11}
Seja $ f(x) = x + \frac{1}{x}$ para $ x > 0$.
\begin{enumerate}
  \item Mostre isso para $0 < u < v $,
        $ f(v) - f(u) = (v - u)\,\frac{uv - 1}{uv}$.
  \item Deduza isso $f$ é estritamente decrescente sobre $\intoc{0}{1}$ e
        estritamente crescente sobre $\intco{1}{+\infty}$.
  \item Conclua isso $ x + \frac1x \geq 2$ para todos $ x > 0$, com igualdade
        apenas em $ x = 1$.
\end{enumerate}
\end{exercise}

\section{Problema: A fábrica de simetria}

\begin{problem}[{Problema de fim de semana --- cada função é par
parte mais uma parte ímpar: regras de paridade, o jogo de transformação e
mínimos ganhos sem cálculo}]
\label{pb:g11:func:1}
Alguns funções são simétricos em espelho, alguns são meia volta
simétrico, a maioria não é nenhum dos dois --- ainda um pequeno milagre algébrico
diz que \emph{todo} \omterm{def:g11:func:function}{função} é, exatamente de uma maneira, a soma
de uma peça simétrica de espelho e uma peça simétrica de meia volta. Isto
problema prova o milagre, joga o jogo da transformação no
curvas de referência (\cref{prop:g11:func:transforms}) e fecha
com o fatoração de diferença técnica de
\cref{exo:g11:func:11}, que vence problemas de minimização um completo
capítulo antes que a derivada chegue.

\medskip
\noindent\textbf{Parte I --- Fluência de paridade.}
\begin{enumerate}
  \item \omterm{def:g11:func:parity}{Par}, \omterm{def:g11:func:parity}{ímpar}, ou nenhum (\cref{def:g11:func:parity})?
        Justifique cada um: $ x^4 - 3x^2$; \ $ x^3 - x $; \ $ x^2 + x $; \
        $\abs x $; \ $\frac1x $.
  \item Indique o significado geométrico de cada veredicto
        (\cref{prop:g11:func:paritysym}): qual simetria faz
        a curva de uma \omterm{def:g11:func:parity}{função par} ter? De um \omterm{def:g11:func:parity}{ímpar} um?
  \item A paridade se multiplica como sinais: prove que o produto de
        dois funções \omterm{def:g11:func:parity}{ímpares} é \omterm{def:g11:func:parity}{par}, e complete todo
        tabuada de multiplicação (\omterm{def:g11:func:parity}{par} $\times $ \omterm{def:g11:func:parity}{par}, \omterm{def:g11:func:parity}{par} $\times $
        \omterm{def:g11:func:parity}{ímpar}, \omterm{def:g11:func:parity}{ímpar} $\times $ \omterm{def:g11:func:parity}{ímpar}).
  \item Consequências rápidas: se $f$ é \omterm{def:g11:func:parity}{ímpar} e definido em $0$,
        o que é $ f(0)$? Se $f$ é \omterm{def:g11:func:parity}{par} com $ f(3) = 7$, o que é
        $ f(-3)$? E mostrar que o único \omterm{def:g11:func:function}{função} ambos \omterm{def:g11:func:parity}{par} e
        \omterm{def:g11:func:parity}{ímpar} é o zero \omterm{def:g11:func:function}{função}.
  \item Quanto de uma \omterm{def:g11:func:function}{função} parA curva de determina todos
        isso? Dê uma \omterm{def:g11:func:function}{função} isso não é nenhum dos dois \omterm{def:g11:func:parity}{par} nem ímpare
        explique o que falha.
\end{enumerate}

\medskip
\noindent\textbf{Parte II --- O teorema da decomposição.}
Por uma \omterm{def:g11:func:function}{função} $f$ definido em um domínio simétrico sobre $0$, definir
\[
E(x) = \frac{f(x) + f(-x)}{2},
\qquad
O(x) = \frac{f(x) - f(-x)}{2} .
\]
\begin{enumerate}[resume]
  \item Calcular $ E $ e $ O $ para $ f(x) = x^3 + x^2 + 1$e
        verifique: $ E $ \omterm{def:g11:func:parity}{par}, $ O $ \omterm{def:g11:func:parity}{ímpar}, $ E + O = f $.
  \item Prove a afirmação geral: para qualquer $f$, a \omterm{def:g11:func:function}{função}
        $ E $ é \omterm{def:g11:func:parity}{par}, $ O $ é \omterm{def:g11:func:parity}{ímpar}, e $ f = E + O $ --- \emph{todo
        \omterm{def:g11:func:function}{função} se divide em um \omterm{def:g11:func:parity}{par} parte e um \omterm{def:g11:func:parity}{ímpar} papel}.
  \item Prove que a divisão é única: se
        $ f = E_1 + O_1 = E_2 + O_2$ com $ E_i $ \omterm{def:g11:func:parity}{par} e $ O_i $
        \omterm{def:g11:func:parity}{ímpar}, use a pergunta 4 em $ E_1 - E_2$.
  \item Decompor $ f(x) = \dfrac{1}{1 - x}$ (para
        $ x \neq \pm 1$): colocar $ E $ e $ O $ acima do comum
        denominador $1 - x^2$ e simplificar.
  \item Um olhar para frente, em uma frase cada: o \omterm{def:g11:func:parity}{par} e \omterm{def:g11:func:parity}{ímpar}
        partes do exponencial \omterm{def:g11:func:function}{função} será o gêmeo
        funções da nota 12 (cosseno e seno hiperbólico); e
        o mesmo instinto --- dividir um objeto em simétrico
        peças e trate cada uma com sua própria simetria --- poderes
        a análise de sons e sinais dos volumes universitários.
        Por que essa divisão \emph{útil}, em geral?
\end{enumerate}

\medskip
\noindent\textbf{Parte III --- O jogo da transformação.}
\begin{enumerate}[resume]
  \item A partir da curva de $ x \mapsto x^2$, descreva
        precisamente as curvas de $(x - 3)^2$, \ $ x^2 + 2$, \
        $(x - 3)^2 + 2$, \ $-x^2$, \ $2x^2$
        (\cref{prop:g11:func:transforms}).
  \item Reduzir $ f(x) = \abs{x - 2} + \abs{x + 2}$ para um
        fórmula por partes nos três \omterm{def:g10:numbers:interval}{intervalos} cortado por $-2$
        e $2$, e descreva sua curva (tem um famoso plano
        inferior).
  \item Quais transformações respeitam a uniformidade? Decida, com
        prova ou contraexemplo: deslocamento horizontal
        $ f(x - 3)$; mudança vertical $ f(x) + 2$; alongamento vertical
        $2 f(x)$ (todos aplicados a um \omterm{def:g11:func:parity}{par} $f$).
  \item Resolver $\abs{x - 2} + \abs{x + 2} = 6$ usando o
        fórmula por partes da questão 12.
  \item Engenharia reversa: uma \omterm{def:g11:quad:parabola}{parábola} tem vértice $(2, -3)$
        e passa por $(0, 1)$. Reconstrua seu equação em
        \omterm{thm:g11:quad:canonical}{forma canônica}.
\end{enumerate}

\medskip
\noindent\textbf{Parte IV --- \omterm{def:g10:functions:extrema}{Mínimos} sem cálculo.}
\begin{enumerate}[resume]
  \item Espelho \cref{exo:g11:func:11} para
        $ g(x) = x + \dfrac4x $ sobre $\intoo{0}{+\infty}$: mostrar
        $ g(v) - g(u) = (v - u)\,\dfrac{uv - 4}{uv}$, localize o
        \omterm{def:g10:functions:extrema}{mínimo}, e dê seu valor.
  \item Aplicar: entre todos os retângulos da área $16$, que tem o
        menor perímetro? (Escrever $ P = 2\left(x +
        \frac{16}{x}\right)$ e use a máquina da pergunta 16.) Um
        conclui o velho amigo.
  \item Determine as variações de
        $ h(x) = \sqrt{x^2 + 1}$ em todos $\R $, combinando um
        argumento de composição em $\intco{0}{+\infty}$ com um
        argumento de paridade para o resto.
  \item O ponto mais próximo: encontre o ponto da curva
        $ y = \sqrt x $ mais próximo de $(3, 0)$. (Minimize o
        \emph{ao quadrado} distância --- uma quadrática em $ x$ --- e
        explique por que minimizar o quadrado é legítimo.)
  \item Finalizando, o tour pela fábrica: as curvas de referência são a matéria-prima
        estoque; transforma as máquinas; paridade e o
        decomposição do controle de qualidade; fatoração de diferença
        a bancada de medição. Uma frase em cada --- e nome
        a ferramenta elétrica que o próximo capítulo apresenta.
\end{enumerate}
\end{problem}
